\section{A joint taxonomy}
\label{sec:taxonomy}

\subsection{Three axes}

We place every named mechanism and every named defense on three axes.

\textbf{Channel} is the route the content takes into the context window. There
are nine: \meth{direct-input}, \meth{RAG}, \meth{tool-output},
\meth{tool-metadata}, \meth{memory}, \meth{skill}, \meth{multi-agent},
\meth{cross-modal} and \meth{supply-chain}.

\textbf{Intent} is whether the content was placed deliberately. There are
three: adversarial, incidental, and both.

\textbf{Consequence} is what the agent then does wrong. There are six:
\meth{goal-hijack}, \meth{reasoning-corruption}, \meth{silent-corruption},
\meth{persistence-backdoor}, \meth{data-exfiltration} and
\meth{resource-abuse}.

The three axes give 162 cells. The axes are chosen so that the two literatures
share them. An attack and an incidental mechanism that enter through the same
channel and cause the same damage occupy the same cell, and differ only in the
intent coordinate. That is what makes the taxonomy a tool for finding missing
experiments rather than only a filing scheme.

\subsection{The cube is mostly empty}

\textbf{98 of the 162 cells (60.5\%) contain no paper from either literature.}
Emptiness is not uniform. Table~\ref{tab:empty} gives the empty-cell count per
channel.

\begin{table}[t]
\caption{Empty cells per channel, out of 18 (3 intents $\times$ 6
consequences).}
\label{tab:empty}
\small
\begin{tabular}{@{}lr@{\hspace{2em}}lr@{}}
\toprule
\textbf{Channel} & \textbf{Empty} & \textbf{Channel} & \textbf{Empty} \\
\midrule
supply-chain  & 15/18 & multi-agent   & 11/18 \\
skill         & 14/18 & direct-input  & 10/18 \\
tool-metadata & 11/18 & tool-output   & 10/18 \\
cross-modal   & 11/18 & RAG           & 9/18 \\
              &       & memory        & 7/18 \\
\bottomrule
\end{tabular}
\end{table}

The two newest channels are the emptiest. Supply-chain and skill are the routes
by which an agent loads third-party capability, and they are the two the field
has studied least.

\subsection{Census of mechanisms}

The corpus names \textbf{221 distinct mechanisms}, or 235 counting those that
exist only inside a benchmark suite. 166 are adversarial attacks from the
security track and 55 are incidental degradation mechanisms from the machine
learning track.

By consequence, \meth{goal-hijack} dominates at 124 of 221 (56.1\%), followed
by \meth{reasoning-corruption} at 46 (20.8\%) and \meth{silent-corruption} at
21 (9.5\%). By channel, \meth{direct-input} (66, 29.9\%) and \meth{tool-output}
(65, 29.4\%) together hold close to 60\% of all named mechanisms, with
\meth{RAG} third at 27 (12.2\%).

Mechanism count and paper count differ by a large factor, and the reason is
structural. 221 of the 984 included papers name a new mechanism and 529 name a
new defense, with 68 papers doing both. \textbf{302 papers name neither.} Those
are benchmarks, surveys, and measurement papers that apply an existing named
mechanism rather than introducing one. A corpus of this kind is mostly papers
that use the field's vocabulary rather than papers that add to it.

\subsection{Census of defenses}

The corpus names \textbf{529 distinct defenses}. The track split is almost
even: 264 from machine learning papers (49.9\%), 260 from security papers
(49.1\%) and 5 from papers that sit in both.

By intervention point, 205 act at reasoning time (38.8\%), 152 at ingestion
(28.7\%), 72 at execution (13.6\%), and 100 specify no single point (18.9\%).

The load-bearing number is what each defense was validated against in its own
paper:

\begin{itemize}
\item \textbf{261 (49.3\%)} were validated against adversarial scenarios only.
\item \textbf{251 (47.4\%)} were validated against incidental scenarios only.
\item \textbf{14 (2.6\%)} were validated against both.
\item 3 report no empirical validation.
\end{itemize}

\textbf{97\% of proposed defenses have an untested cross-threat-model question
hanging over them.} This is the pool the experiments in
Section~\ref{sec:experiments} draw from.

\subsection{Where protection does not follow the problem}

Putting the two censuses on the same axes shows that defenses are not
distributed like the mechanisms they are meant to address
(Table~\ref{tab:bychannel}).

\begin{table}[t]
\caption{Named mechanisms and named defenses per channel. The ratio is
defenses per mechanism. A low ratio marks a channel the field has named
problems for and not answered.}
\label{tab:bychannel}
\small
\begin{tabular}{@{}lrrr@{}}
\toprule
\textbf{Channel} & \textbf{Mechanisms} & \textbf{Defenses} & \textbf{Ratio} \\
\midrule
memory         & 14 &  72 & 5.1 \\
RAG            & 27 &  86 & 3.2 \\
multi-agent    &  9 &  36 & 4.0 \\
tool-output    & 65 & 172 & 2.6 \\
direct-input   & 66 & 127 & 1.9 \\
skill          &  4 &   6 & 1.5 \\
cross-modal    & 18 &  17 & 0.9 \\
tool-metadata  & 12 &  11 & 0.9 \\
\textbf{supply-chain} & \textbf{6} & \textbf{2} & \textbf{0.3} \\
\bottomrule
\end{tabular}
\end{table}

\textbf{Supply-chain is the clearest gap.} It has six named mechanisms and two
named defenses. Cross-modal and tool-metadata are close behind at under one
defense per mechanism. These are the channels through which an agent accepts
third-party components and machine-readable descriptions, which is the
direction the deployment of these systems is moving.

The effort is also split unevenly by threat model. Counting defenses against
the targets in their own track, there are \textbf{4.6 defenses per incidental
failure mode} but only \textbf{1.4 per named attack}. 72.7\% of incidental
mechanisms have at least one defense tested against them; the figure for named
attacks is 43.3\%. The adversarial half of the problem has more named problems
and fewer answers.
